\noindent An on-chain index basket held at fixed weights is usually run as a weighted (geometric-mean) pool that arbitrageurs rebalance. We show that such a pool is a band rebalancer whose band is set by its swap fee $f$ plus the arbitrageurs' execution cost $c$, and derive its loss against a frictionless constant-weight benchmark: $\gamma^{*}(c+\rho)/(f+c+\rho)$, where $\gamma^{*}$ is the basket's excess growth rate and $\rho\approx0.583\,\sigma\sqrt{\Delta t}$ the overshoot between blocks. The formula is exact to first order for two assets and within 1.4\% of a block-by-block simulation across 32 parameter sets; its basket version is within 3\% at fees of 30\bps\ or more. It implies that faster blocks cannot remove the cost term, and that at 10\bps\ of arbitrage cost only a fee of almost 10\% would make a six-token basket cheap to rebalance. A \emph{gated pool} removes the coupling between fee and rebalancing: bonded solvers rebalance it only when its weights drift past a chosen band, at oracle price plus a spread capped on-chain, and retail trades a separate quote anchored to the oracle. Its cost is $sK/(\delta+\din)$. For a six-token crypto basket at a 30\bps\ fee, 12-second blocks and 10\bps\ costs, rebalancing costs 2.53\% of TVL a year in a public pool and 0.10\% in a gated one. We test twelve assumptions. The advantage survives volatility scaling, regimes, trends, mean reversion and jumps, where it widens to 4.10\% against 0.12\%, and disappears only if arbitrage is nearly free ($c<1$\bps) on blocks of 2 seconds or less while solvers charge more than 14--30\bps. The largest risk we found is oracle price selection on rebalancing fills, which would raise the gated pool's cost tenfold; we specify the timestamp rule that closes it. Every derivation and every number is worked out in the appendices.
